\documentclass[12pt,a4paper]{article}
\usepackage[utf8]{inputenc}
\usepackage[T1]{fontenc}
\usepackage[spanish]{babel}
\usepackage{geometry}
\usepackage{xcolor}
\usepackage{listings}
\usepackage{booktabs}
\usepackage{hyperref}
\usepackage{fancyhdr}
\usepackage{titlesec}
\usepackage{array}
\usepackage{longtable}
\usepackage{parskip}
\usepackage{dirtree}
\usepackage{graphicx}
\usepackage{mdframed}
\usepackage{enumitem}

\geometry{top=2.5cm, bottom=2.5cm, left=3cm, right=3cm}

% Colores
\definecolor{codeblue}{RGB}{41, 98, 255}
\definecolor{codered}{RGB}{220, 50, 47}
\definecolor{codegray}{RGB}{128, 128, 128}
\definecolor{codegreen}{RGB}{38, 166, 65}
\definecolor{codepurple}{RGB}{136, 0, 210}
\definecolor{bgcode}{RGB}{245, 247, 250}
\definecolor{accent}{RGB}{0, 112, 243}
\definecolor{headerblue}{RGB}{15, 30, 70}

% Configuración de listings
\lstset{
  backgroundcolor=\color{bgcode},
  basicstyle=\ttfamily\small,
  breaklines=true,
  captionpos=b,
  commentstyle=\color{codegray},
  keywordstyle=\color{codeblue}\bfseries,
  stringstyle=\color{codegreen},
  numberstyle=\tiny\color{codegray},
  numbers=left,
  numbersep=8pt,
  frame=single,
  framesep=4pt,
  rulecolor=\color{codegray!40},
  tabsize=2,
  showstringspaces=false,
  xleftmargin=12pt,
}

% Estilo de encabezados
\titleformat{\section}{\Large\bfseries\color{headerblue}}{}{0em}{\thesection\quad}[\titlerule]
\titleformat{\subsection}{\large\bfseries\color{accent}}{}{0em}{\thesubsection\quad}
\titleformat{\subsubsection}{\normalsize\bfseries\color{black}}{}{0em}{\thesubsubsection\quad}

% Header/Footer
\pagestyle{fancy}
\fancyhf{}
\fancyhead[L]{\small\textcolor{codegray}{SaukCode.cl --- Documentación Técnica}}
\fancyhead[R]{\small\textcolor{codegray}{Vol. I: Arquitectura}}
\fancyfoot[C]{\small\textcolor{codegray}{\thepage}}
\renewcommand{\headrulewidth}{0.4pt}
\renewcommand{\footrulewidth}{0pt}

% Caja de información
\mdfdefinestyle{infobox}{
  linecolor=accent,
  linewidth=1.5pt,
  backgroundcolor=accent!5,
  innertopmargin=10pt,
  innerbottommargin=10pt,
  innerleftmargin=12pt,
  innerrightmargin=12pt,
  roundcorner=4pt,
}

\hypersetup{
  colorlinks=true,
  linkcolor=accent,
  urlcolor=codeblue,
  pdfauthor={SaukCode},
  pdftitle={SaukCode.cl - Documentación Técnica Vol. I},
}

\begin{document}

% Portada
\thispagestyle{empty}
\begin{center}
  \vspace*{3cm}
  {\Huge\bfseries\color{headerblue} SaukCode.cl}\\[0.5cm]
  {\Large\color{accent} Documentación Técnica}\\[0.3cm]
  {\large\color{codegray} Volumen I: Arquitectura y Stack Tecnológico}\\[2cm]
  \rule{10cm}{1.5pt}\\[0.5cm]
  {\normalsize Plataforma web de recopilación de contenido académico}\\[0.3cm]
  {\normalsize Next.js 15 \textbullet\ React 19 \textbullet\ MDX}\\[0.5cm]
  \rule{10cm}{1.5pt}\\[2cm]
  {\small\color{codegray} Generado: Mayo 2026}
\end{center}
\newpage

% Tabla de contenidos
\tableofcontents
\newpage

% =====================================================
\section{Descripción General del Proyecto}
% =====================================================

SaukCode.cl es una plataforma web personal de recopilación y organización de contenido académico. Centraliza material de estudio proveniente de múltiples fuentes:

\begin{itemize}[leftmargin=1.5cm]
  \item \textbf{INACAP}: 25 asignaturas distribuidas en 7 semestres (más pendientes)
  \item \textbf{Bootcamps}: 10+ cursos de dos plataformas (CódigoFacilito y UDD)
  \item \textbf{E-Learning}: 5 plataformas (Udemy, DevTalles, CódigoFacilito, Platzi, Coursera)
  \item \textbf{Idiomas}: 6 idiomas (Inglés, Alemán, Coreano, Mandarín, Japonés, Ruso)
  \item \textbf{LeetCode}: 100+ problemas con soluciones documentadas
  \item \textbf{CodeVault}: Colección de algoritmos y estructuras de datos
\end{itemize}

El sitio sirve como bitácora personal de aprendizaje, con contenido en formato MDX que soporta matemáticas (\LaTeX), bloques de código con resaltado de sintaxis, tablas, quizzes interactivos y más.

\begin{mdframed}[style=infobox]
  \textbf{Estadísticas del proyecto:} 628 archivos MDX de contenido \textbullet\ 25 cursos INACAP \textbullet\ 10+ bootcamps \textbullet\ 100+ problemas LeetCode \textbullet\ 6 idiomas
\end{mdframed}

% =====================================================
\section{Stack Tecnológico}
% =====================================================

\subsection{Framework Principal}

\begin{table}[h]
\centering
\begin{tabular}{lll}
\toprule
\textbf{Tecnología} & \textbf{Versión} & \textbf{Rol} \\
\midrule
Next.js & 15.5.14 & Framework web (App Router) \\
React & 19.0.0 & UI Library \\
JavaScript/JSX & ES2022+ & Lenguaje principal \\
Node.js & 20.x & Runtime de servidor \\
\bottomrule
\end{tabular}
\caption{Framework y runtime}
\end{table}

\subsection{Dependencias de Producción}

\begin{longtable}{lll}
\toprule
\textbf{Paquete} & \textbf{Versión} & \textbf{Propósito} \\
\midrule
\endhead
framer-motion & 12.23.26 & Animaciones y transiciones \\
katex & 0.16.22 & Renderizado de fórmulas matemáticas \\
react-katex & --- & Integración React para KaTeX \\
highlight.js & 11.11.1 & Resaltado de código \\
prismjs & 1.30.0 & Resaltado de código alternativo \\
react-markdown & 10.1.0 & Renderizado de Markdown \\
remark-gfm & --- & GitHub Flavored Markdown \\
remark-math & --- & Soporte de sintaxis matemática \\
rehype-katex & --- & Renderizado KaTeX en MDX \\
@mdx-js/react & --- & Proveedor de componentes MDX \\
@next/mdx & --- & Integración MDX con Next.js \\
lucide-react & --- & Iconos SVG \\
classnames & --- & Utilidad para clases CSS \\
marked & --- & Parser de Markdown \\
\bottomrule
\caption{Dependencias de producción}
\end{longtable}

\subsection{Herramientas de Desarrollo}

\begin{table}[h]
\centering
\begin{tabular}{lll}
\toprule
\textbf{Herramienta} & \textbf{Versión} & \textbf{Rol} \\
\midrule
ESLint & 9.x & Linting de código \\
eslint-config-next & --- & Reglas específicas de Next.js \\
\bottomrule
\end{tabular}
\caption{DevDependencies}
\end{table}

\subsection{Fuentes y Tipografía}

El proyecto utiliza las fuentes de Google:
\begin{itemize}
  \item \textbf{Geist Sans} (\texttt{--font-geist-sans}): fuente principal para texto de interfaz
  \item \textbf{Geist Mono} (\texttt{--font-geist-mono}): fuente monoespaciada para código
\end{itemize}

Estas se cargan mediante \texttt{next/font/google} en el layout raíz para optimización automática.

% =====================================================
\section{Estructura de Directorios}
% =====================================================

La organización del proyecto sigue las convenciones del App Router de Next.js 15:

\begin{lstlisting}[language=bash, caption={Árbol de directorios principal}]
saukcode.cl/
├── src/
│   ├── app/                    # Rutas (Next.js App Router)
│   │   ├── layout.js           # Layout raíz global
│   │   ├── template.js         # Template de página (animaciones)
│   │   ├── page.jsx            # Página home
│   │   ├── globals.css         # Estilos globales
│   │   ├── inacap/             # Sección INACAP
│   │   ├── bootcamp/           # Sección Bootcamp
│   │   ├── elearning/          # Sección E-Learning
│   │   ├── (idiomas)/          # Grupo de rutas: idiomas
│   │   ├── leetcode/           # Sección LeetCode
│   │   ├── codevault/          # Sección CodeVault
│   │   └── portfolio/          # Portafolio personal
│   ├── components/             # Componentes React
│   │   ├── layout/             # Navbar, Breadcrumb, Footer
│   │   ├── common/             # Componentes reutilizables
│   │   ├── ui/                 # Componentes de UI genéricos
│   │   ├── mdx/                # Componentes específicos MDX
│   │   ├── features/           # Componentes por sección
│   │   └── providers/          # Proveedores de contexto
│   ├── data/                   # Metadatos de cursos y secciones
│   │   ├── homeData.js
│   │   ├── inacap/             # 25+ archivos de datos INACAP
│   │   ├── bootcamp/           # 12 archivos de datos
│   │   ├── elearning/          # Datos de plataformas
│   │   ├── idiomas/            # Datos de idiomas
│   │   ├── leetcode/           # Problemas y categorías
│   │   └── codevault/          # Algoritmos
│   ├── content/                # Archivos MDX (628 total)
│   │   ├── inacap/             # Contenido INACAP
│   │   ├── bootcamp/           # Contenido bootcamp
│   │   ├── elearning/          # Contenido e-learning
│   │   ├── idiomas/            # Contenido idiomas
│   │   └── leetcode/           # Soluciones LeetCode
│   ├── hooks/                  # Custom hooks de React
│   ├── utils/                  # Funciones utilitarias
│   ├── shared/                 # Estilos compartidos CSS
│   └── mdx-components.jsx      # Mapeo global de componentes MDX
├── public/                     # Archivos estáticos
├── next.config.mjs             # Configuración de Next.js
├── package.json
└── jsconfig.json               # Alias de rutas (@/*)
\end{lstlisting}

% =====================================================
\section{Sistema de Routing (App Router)}
% =====================================================

El proyecto utiliza el \textbf{App Router} de Next.js 15, con generación estática (\texttt{generateStaticParams}) para todas las páginas de contenido.

\subsection{Convenciones de Rutas}

\begin{table}[h]
\centering
\begin{tabular}{ll}
\toprule
\textbf{Ruta} & \textbf{Descripción} \\
\midrule
\texttt{/} & Página principal \\
\texttt{/inacap} & Lista de asignaturas INACAP \\
\texttt{/inacap/[courseId]} & Detalle de asignatura \\
\texttt{/inacap/[courseId]/[slug]} & Lección específica \\
\texttt{/bootcamp} & Cursos de bootcamp \\
\texttt{/bootcamp/[courseId]/[slug]} & Clase de bootcamp \\
\texttt{/elearning/[courseId]/[slug]} & Clase de e-learning \\
\texttt{/[idioma]/[courseId]/[slug]} & Lección de idioma \\
\texttt{/leetcode} & Lista de problemas \\
\texttt{/leetcode/[problemId]} & Solución de problema \\
\texttt{/codevault} & Colección de algoritmos \\
\texttt{/portfolio} & Portafolio personal \\
\bottomrule
\end{tabular}
\caption{Mapa de rutas del sitio}
\end{table}

\subsection{Patrón de Generación Estática}

Todas las páginas de contenido usan el patrón de generación estática en build time:

\begin{lstlisting}[language=JavaScript, caption={Patrón de generación estática típico}]
// src/app/inacap/[courseId]/[slug]/page.jsx

export const dynamicParams = false; // 404 para rutas no generadas

export async function generateStaticParams() {
  return getAllLessonPaths('inacap'); // Escanea filesystem
}

export default async function LessonPage({ params }) {
  const { courseId, slug } = params;
  const { default: Content } =
    await import(`@/content/inacap/${courseId}/${slug}.mdx`);

  return (
    <CustomMDXProvider>
      <Content />
    </CustomMDXProvider>
  );
}
\end{lstlisting}

Con \texttt{dynamicParams = false}, cualquier ruta no generada en build time retorna un 404 automáticamente.

\subsection{Grupo de Rutas: Idiomas}

Los idiomas usan un \textbf{Route Group} (\texttt{(idiomas)/}) que permite agrupar varias rutas bajo un layout compartido sin afectar la URL:

\begin{lstlisting}[language=bash]
src/app/(idiomas)/
├── chino/[courseId]/[slug]/
├── ingles/[courseId]/[slug]/
├── japones/[courseId]/[slug]/
├── aleman/[courseId]/[slug]/
├── coreano/[courseId]/[slug]/
└── ruso/[courseId]/[slug]/
\end{lstlisting}

% =====================================================
\section{Configuración del Proyecto}
% =====================================================

\subsection{next.config.mjs}

\begin{lstlisting}[language=JavaScript, caption={next.config.mjs}]
import createMDX from '@next/mdx';
import remarkGfm from 'remark-gfm';
import remarkMath from 'remark-math';
import rehypeKatex from 'rehype-katex';

const withMDX = createMDX({
  extension: /\.mdx?$/,
  options: {
    remarkPlugins: [remarkGfm, remarkMath],
    rehypePlugins: [rehypeKatex],
  },
});

const nextConfig = {
  pageExtensions: ['js', 'jsx', 'mdx', 'ts', 'tsx'],
  images: {
    remotePatterns: [{
      protocol: 'https',
      hostname: 'kapowaz.github.io',
      pathname: '/square-flags/flags/**',
    }],
  },
};

export default withMDX(nextConfig);
\end{lstlisting}

Los plugins configurados son:
\begin{itemize}
  \item \textbf{remark-gfm}: Soporte de tablas, listas de tareas y otras extensiones de GitHub Flavored Markdown
  \item \textbf{remark-math}: Detecta bloques de sintaxis matemática (\texttt{\$...\$} y \texttt{\$\$...\$\$})
  \item \textbf{rehype-katex}: Convierte la sintaxis matemática detectada en HTML renderizado con KaTeX
\end{itemize}

\subsection{jsconfig.json --- Alias de Rutas}

\begin{lstlisting}[language=JavaScript, caption={jsconfig.json}]
{
  "compilerOptions": {
    "paths": {
      "@/*": ["./src/*"]
    }
  }
}
\end{lstlisting}

El alias \texttt{@/*} mapea a \texttt{./src/*}, permitiendo imports limpios como:

\begin{lstlisting}[language=JavaScript]
import { usePagination } from '@/hooks/usePagination';
import { buildLessonLink } from '@/utils/linkBuilder';
import styles from '@/shared/layout.module.css';
\end{lstlisting}

\subsection{Scripts NPM}

\begin{table}[h]
\centering
\begin{tabular}{ll}
\toprule
\textbf{Script} & \textbf{Comando} \\
\midrule
\texttt{npm run dev} & Inicia servidor de desarrollo \\
\texttt{npm run build} & Genera build de producción (estático) \\
\texttt{npm run start} & Inicia servidor de producción \\
\texttt{npm run lint} & Ejecuta ESLint \\
\bottomrule
\end{tabular}
\caption{Scripts disponibles}
\end{table}

% =====================================================
\section{Sistema de Temas (Dark/Light Mode)}
% =====================================================

El sistema de temas opera directamente sobre las clases CSS del elemento \texttt{<html>}, sin usar Context API ni variables de estado global.

\subsection{Clases de Tema}

\begin{table}[h]
\centering
\begin{tabular}{ll}
\toprule
\textbf{Clase CSS} & \textbf{Estado} \\
\midrule
\texttt{dark-mode} & Modo oscuro por preferencia del sistema \\
\texttt{dark-mode-override} & Modo oscuro forzado por el usuario \\
\texttt{light-mode-override} & Modo claro forzado por el usuario \\
(sin clase) & Modo claro por defecto del sistema \\
\bottomrule
\end{tabular}
\caption{Clases de control de tema}
\end{table}

\subsection{Flujo de Inicialización}

La lógica de inicialización vive en el \texttt{RootLayout} (con directiva \texttt{'use client'}):

\begin{lstlisting}[language=JavaScript, caption={Inicialización del tema en layout.js}]
useEffect(() => {
  const savedTheme = localStorage.getItem('theme');
  const prefersDark =
    window.matchMedia('(prefers-color-scheme: dark)').matches;

  document.documentElement.classList.remove(
    'dark-mode', 'dark-mode-override', 'light-mode-override'
  );

  if (savedTheme === 'dark') {
    document.documentElement.classList.add('dark-mode-override');
  } else if (savedTheme === 'light') {
    document.documentElement.classList.add('light-mode-override');
  } else if (prefersDark) {
    document.documentElement.classList.add('dark-mode');
  }
}, []);
\end{lstlisting}

\subsection{Persistencia}

La preferencia de tema se almacena en \texttt{localStorage} con clave \texttt{'theme'}. El \texttt{CustomMDXProvider} observa cambios en las clases del \texttt{<html>} mediante un \texttt{MutationObserver} para actualizar estilos inline de componentes MDX.

% =====================================================
\section{Flujo de Datos General}
% =====================================================

El flujo de datos del proyecto sigue un patrón unidireccional de tres capas:

\begin{enumerate}
  \item \textbf{Capa de Datos} (\texttt{src/data/}): Archivos JS que definen metadatos de cursos (títulos, links, estructura de módulos)
  \item \textbf{Capa de Utilidades} (\texttt{src/utils/}): Funciones que parsean los datos y escanean el filesystem para generar rutas
  \item \textbf{Capa de Presentación} (\texttt{src/app/} + \texttt{src/components/}): Páginas Next.js y componentes React que consumen los datos
\end{enumerate}

\begin{lstlisting}[language=bash, caption={Flujo de datos simplificado}]
data/inacap/ti3v53.js          # Define estructura del curso
  └─> utils/contentUtils.js    # getLessonSlugs() escanea filesystem
        └─> app/inacap/[courseId]/[slug]/page.jsx
              └─> generateStaticParams()  # Pre-genera todas las rutas
                    └─> import('@/content/inacap/...')  # Carga MDX
\end{lstlisting}

% =====================================================
\section{Conclusiones de Arquitectura}
% =====================================================

La arquitectura de SaukCode.cl está optimizada para:

\begin{itemize}
  \item \textbf{Rendimiento}: Generación estática completa en build time; no hay peticiones al servidor en tiempo de ejecución para el contenido
  \item \textbf{Escalabilidad de contenido}: Agregar nuevas lecciones solo requiere crear archivos MDX; el sistema las detecta automáticamente
  \item \textbf{Mantenibilidad}: Separación clara entre datos, utilidades y presentación
  \item \textbf{Experiencia de usuario}: Animaciones fluidas con Framer Motion, temas persistentes, navegación sin recarga completa
\end{itemize}

\vspace{1cm}
\begin{center}
\rule{8cm}{0.4pt}\\[0.3cm]
{\small\color{codegray} SaukCode.cl --- Documentación Técnica Vol. I\\
Arquitectura y Stack Tecnológico}
\end{center}

\end{document}
